\documentclass[11pt]{article}
\usepackage{geometry}                % See geometry.pdf to learn the layout options. There are lots.
\geometry{letterpaper}                   % ... or a4paper or a5paper or ... 
%\geometry{landscape}                % Activate for for rotated page geometry
%\usepackage[parfill]{parskip}    % Activate to begin paragraphs with an empty line rather than an indent
\usepackage{graphicx}
\usepackage{amssymb}
\usepackage{epstopdf}
\usepackage[applemac]{inputenc}
\usepackage[german]{babel}
\DeclareGraphicsRule{.tif}{png}{.png}{`convert #1 `dirname #1`/`basename #1 .tif`.png}

\title{Percent to Grade}
\author{P. Bauer}
\date{}                                           % Activate to display a given date or no date

\begin{document}
\maketitle
\section{Einf�hrung}
Schreiben Sie ein Programm {\tt PercentToGrade}, in welchem die Prozente, die bei einer Leistungsbeurteilung erreicht wurden, eingegeben werden k�nnen und das daraufhin die erreichte Note (von ``Sehr Gut'' bis ``Nicht Gen�gend'') ausgibt.

\section{Zuordnung Prozente zu Note}
Die Ermittlung der Note soll aufgrund der folgenden Tabelle erfolgen:

\begin{tabular}{|l|c|}
	\hline
	Prozente & Note \\ \hline
	[88; $\infty$[ & Sehr Gut \\ \hline
	[75; 88[ & Gut \\ \hline
	[63; 75[ & Befriedigend \\ \hline
	[50;  63[ & Gen�gend \\ \hline
	[0; 50[ & Nicht Gen�gend \\ \hline
	sonst & Fehlermeldung \\ \hline
\end{tabular}

\begin{itemize}
	\item Beachten Sie, dass der Prozentwert nicht nur ganzzahlig sein muss.
	\item Falls ein Prozentwert unter 0 eingegeben wird, soll die Fehlermeldung ``Ung�ltige Eingabe'' ausgegeben werden
\end{itemize}

\section{Beispiel}
Folgende Abbildung zeigt, wie sich das Programm mit einer konkreten Eingabe verhalten soll. Beachten Sie, dass Ihre Implementierung genau dieses Verhalten abbilden soll. 

%\begin{figure}
\begin{center}
\includegraphics[scale=.7]{SampleOutput.png}
\end{center}
%\caption{Beispiel}
%\label{fig:insertion_sort}
%\end{figure}



\section{Source Code Qualit�t}
Achten Sie darauf, dass Sie korrekte Einr�ckungen machen und dass die Variablennamen richtig geschrieben werden.

\end{document}  